\documentclass{exam}
\usepackage[utf8]{inputenc}

\begin{document}

\begin{center}
\fbox{\fbox{\parbox{5.5in}{\centering
	This is a multiple choice and free response practice exam. Time yourself for 30 minutes and answer the following questions independently.}}}
	\vspace{5mm} \\
	{\LARGE \textbf{Day 2 Chemistry Homework}}
\end{center}

\vspace{2mm}

\makebox[\textwidth]{Name:\enspace\hrulefill}

\vspace{5mm}

\makebox[\textwidth]{Date:\enspace\hrulefill}

\vspace{2mm}

\begin{questions}
\question What are the 4 types of macromolecules?

\vspace{20mm}

\question What are the 3 types of a protein molecule?

\vspace{20mm}

\question What are the 3 parts of a nucleic acid?

\vspace{20mm}

\question What is the difference between incomplete dominance and codominance? Provide an example of each.

\vspace{20mm}

\question Mendel crossed plants that were homozygous dominant for round yellow peas (RRYY) with plants that were homozygous recessive for wrinkled green peas (rryy). Complete this two-factor cross for the F1  generation.

\vspace{20mm}

\question When Mendel crossed the F1 plants were heterozygous dominant for round yellow peas (RrYy) he was able to produce the F2 generation. Show this cross below.

\vspace{20mm}

\question What is the initial step of primary succession?

\vspace{20mm}

\question What are some examples of when primary succession occurs?

\vspace{20mm}

\question Why is the cell membrane considered a semi-permeable membrane?

\vspace{20mm}

\question List 2 different ways active transport differs from passive transport.

\vspace{20mm}

\question What happens when a cell is in a hypertonic state?

\vspace{20mm}

\question What are the types of endocytosis?

\vspace{20mm}

\question Define the name of the process in which the cell removes waste.

\vspace{20mm}

\question What is one function of lipids?

\begin{oneparchoices}
	\choice Long-term energy storage
	\\
	\choice Phospholipids in the cell membrane
	\\
	\choice Steroids
	\\
	\choice All of the above choices are functions of lipids
\end{oneparchoices}

\question What are the components of nucleic acids?

\begin{oneparchoices}
	\choice Phosphate group, pentose sugar, nitrogenous base
	\\
	\choice Amino group, carboxyl group, R group
	\\
	\choice Glycerol, fatty acids
	\\
	\choice Glucose, fructose, maltose, sucrose
\end{oneparchoices}

\question Name the process which is responsible for the replication of sex/gamete cells.

\begin{oneparchoices}
	\choice Mitosis
	\\
	\choice Meiosis
	\\
	\choice Cytokinesis
	\\
	\choice Metaphase
\end{oneparchoices}

\question Define codominance.

\begin{oneparchoices}
	\choice The dominant and recessive traits of a heterozygous trait gets blended together
	\\
	\choice There are more than 2 genotypes for an organism
	\\
	\choice The dominant and recessive traits of a heterozygous trait both are expressed separately
	\\
	\choice Traits that will be only be expressed when an organism does not have a conflicting dominant trait
\end{oneparchoices}

\question A fire occurs in a forest, burning everything to the ground. What is the first thing that will happen as the forest attempts to regrow?

\begin{oneparchoices}
	\choice Bare rock is broken down by lichens to create soil
	\\
	\choice Large trees immediately begin to grow
	\\
	\choice Small grasses and perennials begin to appear
	\\
	\choice Pioneer species start to grow
\end{oneparchoices}

\question A small, charged molecule attempts to pass through the cell membrane. Determine what happens to the molecule immediately.

\begin{oneparchoices}
	\choice The molecule is not allowed to pass through the cell membrane
	\\
	\choice The molecule is stored in the cell membrane 
	\\
	\choice The molecule freely passes through the cell membrane
	\\
	\choice The charge of the molecule is removed and the molecule passes through the membrane 
\end{oneparchoices}

\question What is an example of passive transport?

\begin{oneparchoices}
	\choice Osmosis
	\\
	\choice Sodium-Potassium Pump
	\\
	\choice Proton Pump
	\\
	\choice Endocytosis
\end{oneparchoices}

\question What is the process by which a cell ingests large particles?

\begin{oneparchoices}
	\choice Pinocytosis
	\\
	\choice Phagocytosis
	\\
	\choice Exocytosis
	\\
	\choice Receptor-mediated endocytosis
\end{oneparchoices}

\question A fire occurs in a forest, burning everything to the ground. What is the first thing that will happen as the forest attempts to regrow?

\begin{oneparchoices}
	\choice Bare rock is broken down by lichens to create soil
	\\
	\choice Large trees immediately begin to grow
	\\
	\choice Small grasses and perennials begin to appear
	\\
	\choice Pioneer species start to grow
\end{oneparchoices}

	\question Which of the following has only the elements carbon, hydrogen, and oxygen in the ratio 1:2:1, respectively?

\begin{oneparchoices}
	\choice Proteins
	\\
	\choice Lipids
	\\
	\choice Carbohydrates
	\\
	\choice Nucleic Acids
\end{oneparchoices}


\question Who is known as the "father of genetics"?

\begin{oneparchoices}
	\choice Richard Feynman
	\\
	\choice Charles Darwin
	\\
	\choice Rosalind Franklin
	\\
	\choice Gregor Mendel
\end{oneparchoices}


\question When is a solution hypertonic with respect to a cell placed in it?

\begin{oneparchoices}
	\choice The solution has a higher solute concentration than the cell
	\\
	\choice The concentration of solute in the cell and in the solution is the same
	\\
	\choice The cell has a higher solute concentration than the solution
	\\
	\choice The solution has a lower solute concentration than the cell
\end{oneparchoices}

\end{questions}

\end{document}